\documentclass[book.tex]{subfiles}
\begin{document}

\chapter{Neural Operators}
\label{chap:neural_operators}
\noindent\rule{11cm}{0.4pt} \\
\textbf{Prerequisites:}  \autoref{chap:pde}, \autoref{chap:numerical_pde}, \autoref{chap:neural_ode}  \\
\textbf{Difficulty Level:} ** \\
\noindent\rule{11cm}{0.4pt}
\vspace{5mm}

%"""
%The classical development of neural networks has primarily focused on learning mappings between finite dimensional Euclidean spaces or finite sets. We propose a generalization of neural networks tailored to learn operators mapping between infinite dimensional function spaces. We formulate the approximation of operators by composition of a class of linear integral operators and nonlinear activation functions, so that the composed operator can approximate complex nonlinear operators. We prove a universal approximation theorem for our construction. The proposed neural operators are resolution-invariant: they share the same network parameters between different discretizations of the underlying function spaces and can be used for zero-shot super-resolutions. Numerically, the proposed models show superior performance compared to existing machine learning based methodologies on Burgers' equation, Darcy flow, and the Navier-Stokes equation, while being several order of magnitude faster compared to conventional PDE solvers.
%"""

Neural network theory and practice has mainly focused on learning functions which operate on finite dimensional inputs. Surprisingly, it turns out that neural methods can be adapted to learn general operators on infinite dimensional function space \cite{li2020fourier}, \cite{li2020neural}. That is, we can model elements of
\begin{align}
    \mathcal{B}(\mathcal{H})
\end{align}
by borrowing ideas from messaging passing neural networks and the classical theory of kernels. Mathematically, these operators are formed by composition of approximate linear integrals and nonlinear activation functions. These composed operators can approximate sophisticated nonlinear operators on function space and in fact satisfy a universal approximation theorem. The representation of these operators is not dependent on a choice of underlying grid. One of the intriging properties of such neural operators is that representations can be trained at one grid resolution and be generalized in a zero-shot fashion to another grid resolution.

The neural operator method can be used to approximately solving partial differential equations by directly learning the solution operator. Variants of the base technique use methods like Fourier transforms to speed up the computation of the kernel. The resulting technique can provide dramatic speedups in solution of large systems of partial differential equations.

\section{Parametric Partial Differential Equations}

A parametric partial differential equation is one whole coefficients are specific by arbitrary functions. We will use the following example in this chapter

\begin{align}
    - \nabla \cdot (a(x) \nabla u(x)) &= f(x),\quad  x \in D \\
    u(x) &= 0,\quad  x \in \partial D
\end{align}



Here $D$ is the domain of the partial differential equation in equation, and $\partial D$ is its boundary. We require that $D$ be a bounded open set. Given data points $\{a_j, u_j\}_{j=1}^{N}$ we want to learn the operator 
\begin{align}
\mathcal{F}: \mathcal{A} \times \Theta \to \mathcal{U}
\end{align}

Where $\mathcal{U}$ and $\mathcal{A}$ are spaces of functions taking values in $\mathbb{R}^d$ for some $d$ (note that $d$ does not have to be the same for both spaces), and $\Theta$ is some space of parameters. The core idea is that solving partial differential equations is slow so we instead directly learn the mapping from data (coefficients and solution points) to potential solutions. We can view this as the task of solving the learning problem
\begin{align}
\min_{\theta \in \Theta} \mathbb{E}_a[\mathcal{L}(\mathcal{F}(a, \theta), \mathcal{F}^*(a, \theta))]
\end{align}
where $\mathcal{F}^*$ is the ideal operator we want to find and $\mathcal{L}$ is a loss functional.

As we have learned in earlier chapters, traditional methods for solving partial differential equations approximate the solution by solving on a mesh. These traditional methods solve one instance of the parameteric PDE with specific coefficients. A neural technique instead learns a solution for a family of parametric PDEs and is fast to evaluate on new instances once the original model has been trained.

\section{Operator Learning}

Most machine learning tasks learn mappings from vectors to vectors. Operator learning seeks to learn a function to function mapping to approximate the true operator solution $\mathcal{F}^*$. To do this, we want to represent functions and operators in a mesh-invariant fashion. That is, our operator must be able to handle inputs that are not on the same mesh points as the training data. We define our neural operators to be of the form

\begin{align}
    u = (K_l \circ \sigma_l \circ \dotsc \circ \sigma_1 \circ K_0)v
\end{align}

Here the $\sigma_i$ are the nonlinear local activations and $K$ are linear operators but non-local. If the original PDE is uniformly elliptic, Green's representation formula implies that we can write solution $u$ in the form.

\begin{align}
    u(x) &= \int_D G_a(x,y)[f(y) + (\Gamma_a u)(y)] dy
\end{align}

\noindent where $G_a$ is a potential function and $\Gamma_a$ is an operator.

To solve more general nonlinear PDEs, we build a neural representation for $u$ that is inspired by Green's representation formula. We introduce network $\kappa_\theta$ and define a kernel operator $\mathcal{K}_a$ as

\begin{align}
    (\mathcal{K}_a u)(x) = \int_D \kappa_\theta(x, y, a(x), a(y)) u(y) dy
\end{align}
Here $x, y \in D$ are spatial locations and $a(x), a(y)$ are parameter values. This base equation isn't sufficient to model the action of operator $\Gamma_a$, so we construct an iterative modification of the base operator $\mathcal{K}_a$. We add iterations for $t=1,\dotsc T$ as follows
\begin{align}
u_{t+1}(x) = \sigma \left ((W  + \mathcal{K}_a)u_t(x) \right )
\end{align}
where $u_0 = a$, $W \in \mathbb{R}$ is a learned weight, and $\sigma$ is a nonlinear activation. Note however that this base formula is a little limiting since we are stuck to one dimension. For this reason, it is useful to define the generalization to higher dimensions. Let
\begin{align}
v(x) \in \mathbb{R}^d
\end{align}
and project the initial $u_0$ to a higher dimensional $v_0$ by a projection matrix. We can then define iterations much as we did before with
\begin{align}
v_0 &= P u_0 \\
v_{t+1}(x) &= \sigma \left ((W  + \mathcal{K}_a)v_t(x) \right ), \qquad t = 1,\dotsc,T\\
u_{T} &= Q v_T
\end{align}
where $P, Q$ are suitable projection matrices. If $x, y \in \mathbb{R}^3$, then we have that the kernel has type
\begin{align}
\kappa_\theta: \mathbb{R}^{2(3+1)} \to \mathbb{R}^{d \times d}
\end{align}
We can slightly generalize the above design to have $P$ and $Q$ be local networks (encoder and decoder) where $P$ lifts the input to a high dimensional channel space in some arbitrary fashion while $Q$ projects back to the original space. The full network can then be written compactly as 
\begin{align}
u = Q(K_a \circ \sigma_l \circ \dotsc \circ \sigma_1 \circ K_a)P v
\end{align}
where $\sigma_l$ encodes the transformation
\begin{align}
\sigma_l \circ \mathcal{K}_a (u) &= \sigma((W + \mathcal{K}_a) u)
\end{align}

This series of approximation works since there exists a universal approximation bound that for any continuous operator defined on a compact domain there exists a two layer neural operator which can model this operator. %(\textcolor{red}{TODO: CITE})

%\begin{align}
%\int_D \kappa_\phi(x, y, a(x), a(y)) v_t(y) \nu_x(dy)
%\end{align}

\section{Fixed Discretization}
Let $D_j$ be some fixed discretization of the input domain $D$. Then the functions $a, u$ become discrete vectors
\begin{align}
a | D_j, u | D_j \in \mathbb{R}^n
\end{align}
for some $n$. The operator $\mathcal{K}_a$ then becomes a $n \times n$ kernel matrix $K_a$ in the one dimensional case. In the higher order case, we have
\begin{align}
(K_a)_{x, y} &= \kappa_\theta(x, y, a(x), a(y)) \in \mathbb{R}^{d \times d}
\end{align}
where $x, y \in D_j$ are points in the discretization, so $K$ is a tensor of shape $(n, n, d, d)$. Computing operator updates in this case becomes simple since all operations can be implemented by matrix multiplications.

\section{Graph Networks for Arbitrary Discretization}

In general though, training data may not belong to one fixed discretization, which means that computing the operator updates becomes infeasible with standard matrix multiplications. In this case, we typically use a graph network approximation

\begin{align}
v_{t+1}(x) &= (\sigma(W + \mathcal{K}_a)v_t)(x) \\
& \approx \sigma \left (W v_{t} + \frac{1}{|N(x)|} \sum_{y \in N(x)} \kappa_\theta (e(x,y)) v_t(y) \right)
\end{align}
where $N(x)$ is the neighborhood of $x$. Given an arbitrary discretization $D_j$, we can set $N(x)$ to be the full discretization. This can become very expensive in general, so a form of neighbor thresholding is usually imposed
\begin{align}
N(x) &= \{y \in D_j | \|y - x\| < B \}
\end{align}
where $B$ is some distance cutoff. More sophisticated methods randomly sample the full discretization and perform a low rank decomposition. 

\section{Fourier Neural Operators}

Fourier neural operators provide one method of facilitating the evaluation of the kernel by performing operator learning steps in Fourier space. Let $\mathcal{F}$ denote the Fourier transform in $d$ dimensions

\begin{align}
(\mathcal{F} f)_j(k) &= \int_D f_j(x) e^{-2i\pi \langle x, k \rangle} dx \\
(\mathcal{F}^{-1} f)_j(k) &= \int_D f_j(k) e^{2i\pi \langle x, k \rangle} d 
\end{align}

The Fourier transform converts the convolution operation into pointwise multiplication so we have the identity that
\begin{align}
(\mathcal{K}_\theta u)(x) &= \kappa_\theta \ast u \\
&= \mathcal{F}^{-1}\left (\mathcal{F}(\kappa_\theta) \cdot \mathcal{F}(u) \right )(x)
\end{align}
We then can learn the kernel directly in Fourier space

\begin{align}
R_\theta &= \mathcal{F}(\kappa_\theta)\\
(\mathcal{K}(\phi) v_t)(x) &= \mathcal{F}^{-1} \left ( R_\theta \cdot (\mathcal{F} u) \right ) (x)
\end{align}
One major advantage of this representation is that the Fourier modes are discrete
\begin{align}
k \in \mathbb{Z}^d
\end{align}
We can truncate the higher order Fourier modes (under the ansatz that high frequency information is generally less useful). Then we only need to track
\begin{align} 
\{R_\theta(k) : k \leq N_{\textrm{cutoff}} \}
\end{align}

For a given discretization $D_j$, we can implement the Fourier operation through a matrix multiplication. We can use the Fast Fourier Transform to speed up this computation. These Fourier layers are natively mesh invariant since the parameterization is performed in Fourier space which doesn't directly depend on the choice of mesh. The learned algorithm can be projected out into a higher mesh than featured in the training data, allowing for easy superresolution from the learned algorithm.

%The Fourier layer has complexity $O(kn)$ for the Fourier transform, the FFT is $O(n\log n)$ and the final part is $O(n)$ the powerful advantage of Fourier layers is that they are resolution invariant and mesh invariant.

%One addition is that you can do superresolution sampling by evaluating the learned operator at a finer mesh than the training data provided.%\textcolor{red}{(TODO: more details needed here)}.

%\section{Computing the Kernel Integral}

%Evaluating the kernel integral we have defined previous can become challenging. A number of numerical and mathematical tricks have been defined to speed up solutions and improve their accuracy. 
%\begin{itemize}
%    \item Graph Neural Operators
%    \item Multipole graph neural operators
%    \item Low rank neural operators
%\end{itemize}

%\section{Navier Stokes Equation}

%\textcolor{red}{TODO: Add more discussion}\\

%\textcolor{red}{TODO: Add discussion of semigroup and Markov neural operators.}\\

%\textcolor{red}{TODO: Discuss Physics Informed Neural Operators (PINO)?}

%\section{Numerical Experiments}

%These methods are potentially orders of magnitude faster than other methods and potentially also work for parabolic and hyperbolic solvers.
\end{document}
